% --- Archon physics formalization source begin ---
% archon:physics
% archon:covers PhyXMiniProblems/problem_phyx_mini_0416.lean
% archon:source-report reports/phyx_mini/problem_phyx_mini_0416.source.json

\chapter{Physics problem phyx\_mini\_0416}
\label{ch:PhyXMiniProblems_problem_phyx_mini_0416}

\paragraph{Problem source.}
\#\# Physical scenario

A monatomic gas fills the left end of the cylinder in figure. At 300 \$	ext\{K\}\$, the gas cylinder length is 10.0 \$	ext\{cm\}\$ and the spring is compressed by 2.0 \$	ext\{cm\}\$.

\#\# Question

How much heat energy must be added to the gas to expand the cylinder length to 16.0 \$	ext\{cm\}\$?

\#\# Answer choices

- A: A: 0 J

- B: B: 18 J

- C: C: 39 J

- D: D: 78 J

\#\# Figure caption (auxiliary; use the image as primary evidence)

The image depicts a cross-sectional diagram of a piston-cylinder setup.

- **Left Section:**
  - A shaded region represents the inside of the cylinder.
  - The area inside the square is labeled as "8.0 cm²".
  - The bottom of this section is labeled with "L", suggesting a length measurement.

- **Right Section:**
  - Contains a spring adjacent to the piston.
  - The spring is set against a labeled area indicating a "Vacuum."
  - Text indicates the spring constant as "2000 N/m."

- **General Description:**
  - There is a clear division between the two sections by a dark grey vertical line representing a piston.
  - The left part inside the cylinder is shaded beige, while the right part is white.
  
The diagram suggests a mechanical setup possibly related to physics or engineering, involving pressure, volume, or force calculations.

\#\# Dataset metadata

Ideal Gases and Kinetic Theory; Multi-Formula Reasoning, Physical Model Grounding Reasoning

\paragraph{Recorded answer/context.}
C: C: 39 J

\paragraph{Figure/image path.}
/root/proposal\_for\_physic/science-mango/phyx\_mini\_run/phyx\_data/test\_image/416.png

\paragraph{Formalization target.}
create a compiling Lean file with sorry bodies at `PhyXMiniProblems/problem\_phyx\_mini\_0416.lean`. The Lean declarations must preserve the physical quantities, dimensions or dimensional roles, figure labels, governing-law hypotheses, and final relation expressed by this problem.
Use Mathlib/Physlib names found through LeanExplore where available. If a physics API is missing, introduce faithful local abstractions rather than scalar placeholder aliases.

\begin{theorem}[Physics formalization target]
\label{thm:physics:phyx_mini_0416:target}
\lean{PhyXMiniProblems.ProblemPhyXMini0416.heat_required_for_spring_loaded_monatomic_gas_expansion}
\uses{def:physics:phyx-mini-0416:phyxminiproblems-problemphyxmini0416-springloadedgassetup, def:physics:phyx-mini-0416:phyxminiproblems-problemphyxmini0416-heataddedinjoules, def:physics:phyx-mini-0416:phyxminiproblems-problemphyxmini0416-matchesproblemscenario, def:physics:phyx-mini-0416:phyxminiproblems-problemphyxmini0416-matchessuppliedfigure, def:physics:phyx-mini-0416:phyxminiproblems-problemphyxmini0416-hasphysicalparameters, def:physics:phyx-mini-0416:phyxminiproblems-problemphyxmini0416-satisfiesquasistaticidealgasspringphysics, def:physics:phyx-mini-0416:phyxminiproblems-problemphyxmini0416-answerchoice, def:physics:phyx-mini-0416:phyxminiproblems-problemphyxmini0416-isclosestdisplayedheatanswer}
The assigned autoformalize agent should translate this physics problem into Lean declarations in the covered file, with theorem and lemma proof bodies written as `by sorry`.
\end{theorem}
\begin{proof}
This is an autoformalization task, not a proof task. Produce statements that can later be proved by the physics prover without weakening the physical model.
\end{proof}

% --- Archon declaration topology begin ---
\section{Declaration topology}
\begin{definition}[Cylinder Length Quantity]
  \label{def:physics:phyx-mini-0416:phyxminiproblems-problemphyxmini0416-cylinderlengthquantity}
  \lean{PhyXMiniProblems.ProblemPhyXMini0416.CylinderLengthQuantity}
  A nonnegative cylinder length carrying the physical dimension of length
\end{definition}
\begin{proof}
  This is the declared model definition of Cylinder Length Quantity.
\end{proof}
\begin{definition}[Gas Volume Quantity]
  \label{def:physics:phyx-mini-0416:phyxminiproblems-problemphyxmini0416-gasvolumequantity}
  \lean{PhyXMiniProblems.ProblemPhyXMini0416.GasVolumeQuantity}
  A nonnegative gas volume carrying the physical dimension length cubed
\end{definition}
\begin{proof}
  This is the declared model definition of Gas Volume Quantity.
\end{proof}
\begin{definition}[Spring Constant Quantity]
  \label{def:physics:phyx-mini-0416:phyxminiproblems-problemphyxmini0416-springconstantquantity}
  \lean{PhyXMiniProblems.ProblemPhyXMini0416.SpringConstantQuantity}
  A nonnegative spring constant. Its dimension is mass per time squared, equivalently force per length (N/m)
\end{definition}
\begin{proof}
  This is the declared model definition of Spring Constant Quantity.
\end{proof}
\begin{definition}[length In Metres]
  \label{def:physics:phyx-mini-0416:phyxminiproblems-problemphyxmini0416-lengthinmetres}
  \lean{PhyXMiniProblems.ProblemPhyXMini0416.lengthInMetres}
  \uses{def:physics:phyx-mini-0416:phyxminiproblems-problemphyxmini0416-cylinderlengthquantity}
  Read a cylinder length or spring compression in metres
\end{definition}
\begin{proof}
  This is the declared model definition of length In Metres.
\end{proof}
\begin{definition}[length In Centimetres]
  \label{def:physics:phyx-mini-0416:phyxminiproblems-problemphyxmini0416-lengthincentimetres}
  \lean{PhyXMiniProblems.ProblemPhyXMini0416.lengthInCentimetres}
  \uses{def:physics:phyx-mini-0416:phyxminiproblems-problemphyxmini0416-cylinderlengthquantity}
  Read a cylinder length or spring compression in centimetres
\end{definition}
\begin{proof}
  This is the declared model definition of length In Centimetres.
\end{proof}
\begin{definition}[area In Square Metres]
  \label{def:physics:phyx-mini-0416:phyxminiproblems-problemphyxmini0416-areainsquaremetres}
  \lean{PhyXMiniProblems.ProblemPhyXMini0416.areaInSquareMetres}
  Read an area in square metres
\end{definition}
\begin{proof}
  This is the declared model definition of area In Square Metres.
\end{proof}
\begin{definition}[area In Square Centimetres]
  \label{def:physics:phyx-mini-0416:phyxminiproblems-problemphyxmini0416-areainsquarecentimetres}
  \lean{PhyXMiniProblems.ProblemPhyXMini0416.areaInSquareCentimetres}
  Read an area in square centimetres, as printed in the figure
\end{definition}
\begin{proof}
  This is the declared model definition of area In Square Centimetres.
\end{proof}
\begin{definition}[volume In Cubic Metres]
  \label{def:physics:phyx-mini-0416:phyxminiproblems-problemphyxmini0416-volumeincubicmetres}
  \lean{PhyXMiniProblems.ProblemPhyXMini0416.volumeInCubicMetres}
  \uses{def:physics:phyx-mini-0416:phyxminiproblems-problemphyxmini0416-gasvolumequantity}
  Read a gas volume in cubic metres
\end{definition}
\begin{proof}
  This is the declared model definition of volume In Cubic Metres.
\end{proof}
\begin{definition}[spring Constant In Newtons Per Metre]
  \label{def:physics:phyx-mini-0416:phyxminiproblems-problemphyxmini0416-springconstantinnewtonspermetre}
  \lean{PhyXMiniProblems.ProblemPhyXMini0416.springConstantInNewtonsPerMetre}
  \uses{def:physics:phyx-mini-0416:phyxminiproblems-problemphyxmini0416-springconstantquantity}
  Read a spring constant in newtons per metre
\end{definition}
\begin{proof}
  This is the declared model definition of spring Constant In Newtons Per Metre.
\end{proof}
\begin{definition}[pressure In Pascals]
  \label{def:physics:phyx-mini-0416:phyxminiproblems-problemphyxmini0416-pressureinpascals}
  \lean{PhyXMiniProblems.ProblemPhyXMini0416.pressureInPascals}
  Read a physical pressure in pascals
\end{definition}
\begin{proof}
  This is the declared model definition of pressure In Pascals.
\end{proof}
\begin{definition}[temperature In Kelvin]
  \label{def:physics:phyx-mini-0416:phyxminiproblems-problemphyxmini0416-temperatureinkelvin}
  \lean{PhyXMiniProblems.ProblemPhyXMini0416.temperatureInKelvin}
  Read an absolute temperature in kelvin, accounting for the zero-preserving unit in which a Physlib Temperature value is stored
\end{definition}
\begin{proof}
  This is the declared model definition of temperature In Kelvin.
\end{proof}
\begin{definition}[energy In Joules]
  \label{def:physics:phyx-mini-0416:phyxminiproblems-problemphyxmini0416-energyinjoules}
  \lean{PhyXMiniProblems.ProblemPhyXMini0416.energyInJoules}
  Read heat, internal energy, or work in joules
\end{definition}
\begin{proof}
  This is the declared model definition of energy In Joules.
\end{proof}
\begin{definition}[Process Stage]
  \label{def:physics:phyx-mini-0416:phyxminiproblems-problemphyxmini0416-processstage}
  \lean{PhyXMiniProblems.ProblemPhyXMini0416.ProcessStage}
  The initial and requested final equilibrium states of the gas
\end{definition}
\begin{proof}
  This is the declared model definition of Process Stage.
\end{proof}
\begin{definition}[Chamber Side]
  \label{def:physics:phyx-mini-0416:phyxminiproblems-problemphyxmini0416-chamberside}
  \lean{PhyXMiniProblems.ProblemPhyXMini0416.ChamberSide}
  The left and right regions separated by the piston in the figure
\end{definition}
\begin{proof}
  This is the declared model definition of Chamber Side.
\end{proof}
\begin{definition}[Chamber Contents]
  \label{def:physics:phyx-mini-0416:phyxminiproblems-problemphyxmini0416-chambercontents}
  \lean{PhyXMiniProblems.ProblemPhyXMini0416.ChamberContents}
  Material shown in each region of the supplied figure
\end{definition}
\begin{proof}
  This is the declared model definition of Chamber Contents.
\end{proof}
\begin{definition}[Gas Kind]
  \label{def:physics:phyx-mini-0416:phyxminiproblems-problemphyxmini0416-gaskind}
  \lean{PhyXMiniProblems.ProblemPhyXMini0416.GasKind}
  The gas model specified in the problem statement
\end{definition}
\begin{proof}
  This is the declared model definition of Gas Kind.
\end{proof}
\begin{definition}[Process Kind]
  \label{def:physics:phyx-mini-0416:phyxminiproblems-problemphyxmini0416-processkind}
  \lean{PhyXMiniProblems.ProblemPhyXMini0416.ProcessKind}
  The process idealization needed for equilibrium force balance at each endpoint
\end{definition}
\begin{proof}
  This is the declared model definition of Process Kind.
\end{proof}
\begin{definition}[Piston Cylinder Figure]
  \label{def:physics:phyx-mini-0416:phyxminiproblems-problemphyxmini0416-pistoncylinderfigure}
  \lean{PhyXMiniProblems.ProblemPhyXMini0416.PistonCylinderFigure}
  \uses{def:physics:phyx-mini-0416:phyxminiproblems-problemphyxmini0416-springconstantquantity, def:physics:phyx-mini-0416:phyxminiproblems-problemphyxmini0416-chamberside, def:physics:phyx-mini-0416:phyxminiproblems-problemphyxmini0416-chambercontents}
  The information represented in the primary image. The area and spring constant remain physical quantities; the Boolean fields record the labelled geometry rather than introducing any thermodynamic conclusion
\end{definition}
\begin{proof}
  This is the declared model definition of Piston Cylinder Figure.
\end{proof}
\begin{definition}[Monatomic Gas State]
  \label{def:physics:phyx-mini-0416:phyxminiproblems-problemphyxmini0416-monatomicgasstate}
  \lean{PhyXMiniProblems.ProblemPhyXMini0416.MonatomicGasState}
  \uses{def:physics:phyx-mini-0416:phyxminiproblems-problemphyxmini0416-cylinderlengthquantity, def:physics:phyx-mini-0416:phyxminiproblems-problemphyxmini0416-gasvolumequantity}
  A gas state stores independent physical observables. In particular, volume, pressure, temperature, spring compression, and internal energy are not defined from the answer sought; the governing relations between them are assumptions below
\end{definition}
\begin{proof}
  This is the declared model definition of Monatomic Gas State.
\end{proof}
\begin{definition}[Spring Loaded Gas Setup]
  \label{def:physics:phyx-mini-0416:phyxminiproblems-problemphyxmini0416-springloadedgassetup}
  \lean{PhyXMiniProblems.ProblemPhyXMini0416.SpringLoadedGasSetup}
  \uses{def:physics:phyx-mini-0416:phyxminiproblems-problemphyxmini0416-processstage, def:physics:phyx-mini-0416:phyxminiproblems-problemphyxmini0416-gaskind, def:physics:phyx-mini-0416:phyxminiproblems-problemphyxmini0416-processkind, def:physics:phyx-mini-0416:phyxminiproblems-problemphyxmini0416-pistoncylinderfigure, def:physics:phyx-mini-0416:phyxminiproblems-problemphyxmini0416-monatomicgasstate}
  The closed gas-piston-spring experiment. amountTimesGasConstant is the fixed nR coefficient of this particular gas sample, calibrated in joules per kelvin. Heat and work are independent energy observables
\end{definition}
\begin{proof}
  This is the declared model definition of Spring Loaded Gas Setup.
\end{proof}
\begin{definition}[heat Added In Joules]
  \label{def:physics:phyx-mini-0416:phyxminiproblems-problemphyxmini0416-heataddedinjoules}
  \lean{PhyXMiniProblems.ProblemPhyXMini0416.heatAddedInJoules}
  \uses{def:physics:phyx-mini-0416:phyxminiproblems-problemphyxmini0416-energyinjoules, def:physics:phyx-mini-0416:phyxminiproblems-problemphyxmini0416-springloadedgassetup}
  The joule readout of the heat added to the gas
\end{definition}
\begin{proof}
  This is the declared model definition of heat Added In Joules.
\end{proof}
\begin{definition}[Matches Problem Scenario]
  \label{def:physics:phyx-mini-0416:phyxminiproblems-problemphyxmini0416-matchesproblemscenario}
  \lean{PhyXMiniProblems.ProblemPhyXMini0416.MatchesProblemScenario}
  \uses{def:physics:phyx-mini-0416:phyxminiproblems-problemphyxmini0416-lengthincentimetres, def:physics:phyx-mini-0416:phyxminiproblems-problemphyxmini0416-temperatureinkelvin, def:physics:phyx-mini-0416:phyxminiproblems-problemphyxmini0416-springloadedgassetup}
  The qualitative model and the numerical endpoint data stated in the prose. The final spring compression, final pressure, final temperature, work, and heat are deliberately absent
\end{definition}
\begin{proof}
  This is the declared model definition of Matches Problem Scenario.
\end{proof}
\begin{definition}[Matches Supplied Figure]
  \label{def:physics:phyx-mini-0416:phyxminiproblems-problemphyxmini0416-matchessuppliedfigure}
  \lean{PhyXMiniProblems.ProblemPhyXMini0416.MatchesSuppliedFigure}
  \uses{def:physics:phyx-mini-0416:phyxminiproblems-problemphyxmini0416-areainsquarecentimetres, def:physics:phyx-mini-0416:phyxminiproblems-problemphyxmini0416-springconstantinnewtonspermetre, def:physics:phyx-mini-0416:phyxminiproblems-problemphyxmini0416-springloadedgassetup}
  Facts read directly from the supplied raster: gas on the left, vacuum and the spring on the right, the piston and spring geometry, the L label, area 8.0 cm², and spring constant 2000 N/m. No heat value occurs here
\end{definition}
\begin{proof}
  This is the declared model definition of Matches Supplied Figure.
\end{proof}
\begin{definition}[Has Physical Parameters]
  \label{def:physics:phyx-mini-0416:phyxminiproblems-problemphyxmini0416-hasphysicalparameters}
  \lean{PhyXMiniProblems.ProblemPhyXMini0416.HasPhysicalParameters}
  \uses{def:physics:phyx-mini-0416:phyxminiproblems-problemphyxmini0416-lengthinmetres, def:physics:phyx-mini-0416:phyxminiproblems-problemphyxmini0416-areainsquaremetres, def:physics:phyx-mini-0416:phyxminiproblems-problemphyxmini0416-volumeincubicmetres, def:physics:phyx-mini-0416:phyxminiproblems-problemphyxmini0416-springconstantinnewtonspermetre, def:physics:phyx-mini-0416:phyxminiproblems-problemphyxmini0416-pressureinpascals, def:physics:phyx-mini-0416:phyxminiproblems-problemphyxmini0416-temperatureinkelvin, def:physics:phyx-mini-0416:phyxminiproblems-problemphyxmini0416-springloadedgassetup}
  Positivity and nondegeneracy conditions selecting physical states
\end{definition}
\begin{proof}
  This is the declared model definition of Has Physical Parameters.
\end{proof}
\begin{definition}[Satisfies Quasistatic Ideal Gas Spring Physics]
  \label{def:physics:phyx-mini-0416:phyxminiproblems-problemphyxmini0416-satisfiesquasistaticidealgasspringphysics}
  \lean{PhyXMiniProblems.ProblemPhyXMini0416.SatisfiesQuasistaticIdealGasSpringPhysics}
  \uses{def:physics:phyx-mini-0416:phyxminiproblems-problemphyxmini0416-lengthinmetres, def:physics:phyx-mini-0416:phyxminiproblems-problemphyxmini0416-areainsquaremetres, def:physics:phyx-mini-0416:phyxminiproblems-problemphyxmini0416-volumeincubicmetres, def:physics:phyx-mini-0416:phyxminiproblems-problemphyxmini0416-springconstantinnewtonspermetre, def:physics:phyx-mini-0416:phyxminiproblems-problemphyxmini0416-pressureinpascals, def:physics:phyx-mini-0416:phyxminiproblems-problemphyxmini0416-temperatureinkelvin, def:physics:phyx-mini-0416:phyxminiproblems-problemphyxmini0416-energyinjoules, def:physics:phyx-mini-0416:phyxminiproblems-problemphyxmini0416-springloadedgassetup, def:physics:phyx-mini-0416:phyxminiproblems-problemphyxmini0416-heataddedinjoules}
  The governing endpoint and energy-balance laws in coherent SI readouts: cylinder geometry gives V = A L; piston displacement changes spring compression by the same amount; the sealed ideal gas obeys PV = nRT; quasistatic force balance against vacuum gives PA = kx; a monatomic ideal gas has ΔU = (3/2)nRΔT; work done by the gas is the increase in spring energy; and the first law gives Q = ΔU + W\_by. These are general physical relations. None states the requested heat or an answer-choice value
\end{definition}
\begin{proof}
  This is the declared model definition of Satisfies Quasistatic Ideal Gas Spring Physics.
\end{proof}
\begin{definition}[Answer Choice]
  \label{def:physics:phyx-mini-0416:phyxminiproblems-problemphyxmini0416-answerchoice}
  \lean{PhyXMiniProblems.ProblemPhyXMini0416.AnswerChoice}
  Labels of the four answer choices supplied with the problem
\end{definition}
\begin{proof}
  This is the declared model definition of Answer Choice.
\end{proof}
\begin{definition}[displayed Heat In Joules]
  \label{def:physics:phyx-mini-0416:phyxminiproblems-problemphyxmini0416-displayedheatinjoules}
  \lean{PhyXMiniProblems.ProblemPhyXMini0416.displayedHeatInJoules}
  \uses{def:physics:phyx-mini-0416:phyxminiproblems-problemphyxmini0416-answerchoice}
  Heat in joules printed beside each answer label
\end{definition}
\begin{proof}
  This is the declared model definition of displayed Heat In Joules.
\end{proof}
\begin{definition}[recorded Dataset Answer]
  \label{def:physics:phyx-mini-0416:phyxminiproblems-problemphyxmini0416-recordeddatasetanswer}
  \lean{PhyXMiniProblems.ProblemPhyXMini0416.recordedDatasetAnswer}
  \uses{def:physics:phyx-mini-0416:phyxminiproblems-problemphyxmini0416-answerchoice}
  The source dataset records choice C; this is metadata, not a premise
\end{definition}
\begin{proof}
  This is the declared model definition of recorded Dataset Answer.
\end{proof}
\begin{definition}[Is Closest Displayed Heat Answer]
  \label{def:physics:phyx-mini-0416:phyxminiproblems-problemphyxmini0416-isclosestdisplayedheatanswer}
  \lean{PhyXMiniProblems.ProblemPhyXMini0416.IsClosestDisplayedHeatAnswer}
  \uses{def:physics:phyx-mini-0416:phyxminiproblems-problemphyxmini0416-answerchoice, def:physics:phyx-mini-0416:phyxminiproblems-problemphyxmini0416-displayedheatinjoules}
  A choice is at least as close to a heat readout as every displayed choice
\end{definition}
\begin{proof}
  This is the declared model definition of Is Closest Displayed Heat Answer.
\end{proof}
% --- Archon declaration topology end ---
% --- Archon physics formalization source end ---
